\documentclass[11pt]{article}
\usepackage[margin=0.9in]{geometry}
\usepackage{amsmath,amssymb,amsthm,booktabs,mathrsfs}
\usepackage[T1]{fontenc}
\usepackage{lmodern,microtype}
\usepackage[colorlinks=true,linkcolor=blue,urlcolor=blue]{hyperref}
\newtheorem{theorem}{Theorem}[section]
\newtheorem{lemma}[theorem]{Lemma}
\newtheorem{proposition}[theorem]{Proposition}
\newtheorem{corollary}[theorem]{Corollary}
\theoremstyle{remark}
\newtheorem{remark}[theorem]{Remark}
\newcommand{\norm}[1]{\left\|#1\right\|}
\title{Toward a four-dimensional Yang--Mills construction\\
\large Notebook 30: quantum descent and renormalized block stability}
\author{Research working document}
\date{17 October 2026}
\begin{document}
\maketitle
\begin{abstract}
We continue the existence argument by constructing the local quantum
update on the complete interaction quotient of Notebook 29.
Actual finite minimum comparisons supply connecting gauges that,
in the infinite local limit, approach one constant exponentially.
Compact gauge approximation then proves that classical substitution
preserves local gauge invariance for every invariant input class.
A norm-preserving periodic construction identifies finite normalized
interactions independently of their carrier representation, including
wrapping supports and exact finite Ward identities.

The actual strict Gaussian expansion is continuous in its local
analytic data. Its proof uses support incidence, strict marked-order
tails, an explicit positive inherited-cardinality margin and finite
dimensional dominated convergence. These estimates define the
homogeneous local quantum map on the invariant quotient and bound
the input dependence of its quantum remainder by
\(C(gT/b)\). Finite translation averaging identifies limits of varying
finite inputs with the same map, giving compatibility for every fixed
finite history inside the common budgets.

We then examine the renormalized recurrence using the actual classical
reference actions and the finite endpoint ratio \(R^q\).
Exact reference transport cancels intermediate classical projections.
A fixed outer field domain and vanishing total linear jet give
an \(O(b^2)\) zero-input quantum force on the inner radius \(b\).
The extracted activities therefore have an \(O(b^2)\) invariant
envelope under a finite block of sufficiently weak steps, with
all intermediate inputs in the common domain. Every quantum
correction and vacuum factor remains present. These estimates concern
the homogeneous local sector relative to one fixed joint classical
system. The running
coupling must obey its actual extraction equation; a prescribed
sequence of small parameters does not by itself furnish that
trajectory. Mixed and large-field histories, sourced bounds and
the continuum observable theory remain necessary for existence.
\end{abstract}

The objective is a nontrivial four-dimensional Yang--Mills quantum
theory. All results here use the finite constructions of Notebooks
15--29 with their stated common domains and actual retained-link
windows. The proofs below add explicit bridges between finite
normalized integrals, infinite analytic interaction classes and
renormalized activity estimates. They do not assume that equality
of infinite compact-field totals automatically identifies a finite
integral, nor that a final norm bound alone proves continuity.



% No preamble. Prefix csgq-. Fixed joint homogeneous background system.
\section{Summable connecting gauges and the invariant quotient}
\label{csgq-section}

This section closes the gauge-invariance interface for classical
substitution on the analytic interaction quotient. Its finite-regulator
inputs are the actual ordinary-minimum comparison and equality theorem,
the ordinary-to-corrected transfer, the constraint-preserving flat-slice
flow, and the coherent pre-gauge support estimates. It does not assume a
weighted estimate for the nonlinear gauge-fixing inverse.

\subsection{Joint limits, domains and the statement}
\label{csgq-input-section}

Fix a preceding homogeneous refinement $M$, a finite allowed endpoint
ratio $L$, and a single joint subsequence of compatible cubic regulators.
Let
\begin{equation}
 e^{ia_{L,N}(w)}=P_{M,L,N}(e^{iw})
       \longrightarrow e^{ia_L(w)}
 \label{csgq-background-limit}
\end{equation}
locally at every current link, locally uniformly in the finitely many
variables of each compactly supported final field $w$. The current
lattice has spacing $1/L$ in final physical units. Assume the established
coherent support bound
\begin{equation}
 \|a_L(w)|_i\|_\infty
 \le \rho e^{-\kappa\operatorname{dist}(i,K)},\qquad
 \rho=\frac{C_A b}{L},
 \qquad \|w\|_\infty<b,\quad \operatorname{supp}w\subset K.
 \label{csgq-background-decay}
\end{equation}
Here $K$ is a finite set of final cells. The same assertion holds for
every other compact datum in the ball. Joint analytic compactness and
the finite-regulator support estimates supply this input; unrelated
pointwise choices of subsequences do not replace it.

Use a unitarily invariant matrix norm for group elements. Its comparison
with the scalar-coordinate maximum norm defining the activity polydisc
has fixed color-dependent constants. In particular, constant conjugation
is not silently treated as an isometry of that coordinate maximum norm.
All such comparison constants are included below. The dimension is
$d\ge2$, with $d=4$ in the application. Constants in summability estimates
may depend on fixed $L$ and on a finite compact support, but not on
preceding refinement or approximating volume.

Let $\mathcal Q_{\rm inv}$ be the quotient of anchored analytic
interactions by equality of their vacuum-normalized totals on compact
fields. Its invariant totals are invariant under constant global
adjoint transformations and under compact local gauge paths that stay
inside the analytic field ball. An identity for arbitrary disconnected
endpoint gauges is not part of the conclusion. The summable-field
extension of the quotient gives, for $0<\rho_1<b$,
\begin{equation}
 |\Phi_E(A)-\Phi_E(B)|
 \le \frac{\|[E]\|_{\mathcal Q}}{b-\rho_1}
                 \|A-B\|_{\ell^1_{\rm cell}},
 \qquad \|A\|_\infty,\|B\|_\infty\le\rho_1.
 \label{csgq-summable-continuity}
\end{equation}
The normalized total is absolutely convergent on these summable fields,
and its value depends only on the class.

\begin{theorem}[Universal local gauge identity for a fixed joint map]
\label{csgq-main}
Under these inputs, the scalar substitution
\begin{equation}
 [E]\longmapsto
       \bigl(w\longmapsto \Phi_E(a_L(w))\bigr)
 \label{csgq-substitution}
\end{equation}
preserves local gauge invariance for every invariant input class on a
common smaller germ. If
\begin{equation}
             C_g C_A/L<1/2,
 \label{csgq-margin-condition}
\end{equation}
where $C_g$ is the fixed group/chart constant constructed below, the
compact connecting paths have a common strict margin in the original
radius-$b$ field ball. The local analytic identity then holds throughout
the connected common coarse domain:
\begin{equation}
             \Phi_E(a_L(w^v))=\Phi_E(a_L(w))
 \label{csgq-gauge-identity}
\end{equation}
for the allowed compact local gauge transformations.

The connecting gauges and domain choices are independent of $E$.
Consequently, when the separately established connected representation
and boundedness estimates for classical substitution apply,
\eqref{csgq-substitution} is an action on the whole invariant quotient
for this fixed joint background system.
\end{theorem}

\subsection{The small gauge supplied by the actual equality proof}
\label{csgq-finite-section}

We record the quantitative consequence of the ordinary-minimum equality
proof. For an ordinary-feasible fine field $U=e^{i\eta A}$, the first
transfer uses $g_1=e^{iH_b\theta}$. Its induced-gauge implicit equation
and the cumulative gauge $b_j(U)$ obey
\begin{equation}
 \|\theta\|_\infty\le2\|\omega\|_\infty,\qquad
 \omega=\log b_j(U)/i,\qquad
 \|\omega\|_\infty\le C\|A\|_{C^{1,\alpha}}.
 \label{csgq-first-transfer}
\end{equation}
The proved bound $H_b:\ell^\infty\to C^{2,\alpha}$ controls this gauge
uniformly in mesh and volume. The constraint-preserving flat-slice flow
has scalar velocity
\begin{equation}
 \lambda(t)=\mathcal L_{A(t)}A(t),\qquad
 \|\lambda(t)\|_{C^{2,\alpha}}
 \le C e^{-t}\|R_0D^T A(0)\|_{C^\alpha}
 \le C e^{-t}\|A(0)\|_{C^{1,\alpha}}.
 \label{csgq-flow}
\end{equation}
Integrating its actual unitary gauge ODE therefore gives
\begin{equation}
 \|g_2-I\|_\infty
 \le C\int_0^\infty\|\lambda(t)\|_\infty\,dt
 \le C\|A(0)\|_{C^{1,\alpha}}.
 \label{csgq-integrated-gauge}
\end{equation}
Finally the selected ordinary representative is obtained from the
strict corrected minimizer by
$h_*=e^{iE_{\rm pt}\omega_*}$, where
$\|\omega_*\|_\infty\le C\|B\|_\infty$ and
$E_{\rm pt}:\ell^\infty\to C^{2,\alpha}$ is the fixed smooth
point-value extension.

If equality holds in the two minimum comparisons, strict convexity on
the fixed slice identifies the intervening field with that corrected
minimizer. The product of the actual transfer gauges consequently maps
$U$ to the selected ordinary minimizer and satisfies
\begin{equation}
 \|g_{\rm eq}-I\|_\infty
       \le C\bigl(\|A\|_{C^{1,\alpha}}+\|B\|_\infty\bigr).
 \label{csgq-equality-gauge}
\end{equation}
The constants are uniform on the established smaller comparison balls.
This is stronger than equality of action values alone. It uses no
inverse massless gauge Laplacian and no unproved spatial locality of
$\mathcal L_A$.

For a compact coarse gauge $v=e^{i\theta}$, extend its prescribed values
to the original fine lattice by
$\widehat v=e^{iE_{\rm pt}\theta}$. This extension is compact in physical
space and has $C^{2,\alpha}$ norm bounded by $C\|\theta\|_\infty$,
uniformly in the fine mesh. The fields
$U_{\rm ord}(w)^{\widehat v}$ and $U_{\rm ord}(w^v)$ obey the same
ordinary constraint. The two minimum comparisons, one using the
inverse gauge, give equality. Equation \eqref{csgq-equality-gauge}
therefore provides an actual fine gauge $g_N$ with
\begin{equation}
 U_{\rm ord,N}(w^v)=U_{\rm ord,N}(w)^{g_N},\qquad
 \|g_N-I\|_\infty\le C(b+\|\theta\|_\infty)\le\delta.
 \label{csgq-finite-orbit}
\end{equation}
Choose the data and local gauge neighborhoods so that $\delta$ is a
fixed small number. Every competitor has complete fine-potential norm
at most $C(b+\|\theta\|_\infty)$ and hence lies in the actual comparison
class. Ordinary covariance of the prefix gives
\begin{equation}
 e^{ia_{L,N}(w^v)}
       =\bigl(e^{ia_{L,N}(w)}\bigr)^{g_N|_{\rm current}}.
 \label{csgq-prefix-orbit}
\end{equation}

\subsection{An infinite connecting gauge and its tail}
\label{csgq-tail-section}

Fix real compact $w,\theta$ in these neighborhoods. The endpoint
configurations already converge along the specified joint subsequence.
Compactness of the group and a diagonal subsequence over current sites
give a limiting gauge for this pair. Taking a further subsequence of
gauges changes neither specified endpoint limit. With
$U=e^{ia_L(w)}$ and $V=e^{ia_L(w^v)}$, local passage gives
\begin{equation}
 V_{xy}=g(x)U_{xy}g(y)^{-1},\qquad
                 \sup_x\|g(x)-I\|\le\delta.
 \label{csgq-infinite-orbit}
\end{equation}
No simultaneous or holomorphic choice of $g$ for all pairs is required.
Both endpoints have exponential decay because $w,w^v$ are compact.

\begin{lemma}[One constant at infinity]
\label{csgq-constant-lemma}
Let compact-group link fields $U,V$ on $\mathbb Z^d$, $d\ge2$, satisfy
$V=U^g$. Suppose their positive bonds based at $x$ obey
\begin{equation}
       \|U_b-I\|+\|V_b-I\|\le C_K e^{-\mu|x|_\infty}.
 \label{csgq-endpoint-tail}
\end{equation}
Then there is a single $c\in G$ such that
\begin{equation}
 \|g(x)-c\|\le C'_K(1+|x|_\infty)e^{-\mu|x|_\infty}.
 \label{csgq-gauge-tail}
\end{equation}
Thus $c^{-1}g-I$ is summable. If $\|g-I\|_\infty\le\delta$, its
supremum is at most $2\delta$.
\end{lemma}
\begin{proof}
The edge equation implies $g(y)=V_{xy}^{-1}g(x)U_{xy}$.
Unitary invariance of the matrix norm gives
\begin{equation}
       \|g(y)-g(x)\|\le\|U_{xy}-I\|+\|V_{xy}-I\|.
 \label{csgq-edge-bound}
\end{equation}
The ray $g(ne_1)$ is Cauchy, with tail $O(e^{-\mu n})$, and
its limit $c$ belongs to $G$. For $n=|x|_\infty$, the lattice boundary
of $[-n,n]^d$ connects $x$ to $ne_1$ by a path of length at most
$C_dn$. A boundary coordinate can be held at $\pm n$ while the other
coordinates are adjusted; a second boundary coordinate permits changing
its sign. This uses $d\ge2$. Each positive lower endpoint has norm at
least $n-1$. Summing \eqref{csgq-edge-bound} along the path gives
$C_dnC_Ke^{-\mu(n-1)}$. Combining this with the ray tail proves
\eqref{csgq-gauge-tail}. Polynomial shell growth proves summability.
Taking the limit in the small-supremum bound and using
$\|c^{-1}g-I\|=\|g-c\|$ gives the last assertion.
\end{proof}

For the actual fields, \eqref{csgq-background-decay} implies
\eqref{csgq-endpoint-tail} with $\mu$ a positive multiple of
$\kappa/L$, after enclosing their finite final supports in a cube.
Support-dependent constants are harmless for equality of evaluations.
In one dimension the two ends can have different constants; the
restriction $d\ge2$ is substantive.

\subsection{Compact gauge paths with a common strict margin}
\label{csgq-path-section}

Put $h=c^{-1}g$. Then $h-I\in\ell^1$ and
\begin{equation}
                         U^h=c^{-1}Vc.
 \label{csgq-normalized-gauge}
\end{equation}
For sufficiently small $\delta$ choose the common Hermitian logarithm
$h(x)=e^{i\lambda(x)}$. It has small supremum and belongs to $\ell^1$.
Let $\lambda_n$ equal $\lambda$ inside a current cube of radius $n$
and vanish outside, and set $h_n=e^{i\lambda_n}$. These gauges are
compact and
\begin{equation}
 \|h_n-h\|_{\ell^1}\longrightarrow0,\qquad
 \|U^{h_n}_b-U^h_b\|
       \le\|h_n(x)-h(x)\|+\|h_n(y)-h(y)\|.
 \label{csgq-compact-gauge-limit}
\end{equation}
Hence the transformed links, and their potentials on a common strict
logarithm chart, converge in cell $\ell^1$. A smooth physical cutoff is
not needed on this fixed current lattice: the truncation jump itself
decays exponentially.

The paths, not just their endpoints, stay in the analytic domain.
The edge equation and the Lipschitz bound for the small logarithm give
\begin{equation}
 \|\lambda(y)-\lambda(x)\|
 \le C\bigl(\|a_L(w)_b\|+\|a_L(w^v)_b\|\bigr)\le C\rho.
 \label{csgq-log-gradient}
\end{equation}
For all sufficiently large $n$, the exterior truncation jump of
$\lambda_n$ is also at most $\rho$. Duhamel's formula for real
Hermitian matrices yields, for $0\le t\le1$,
\[
 \|e^{it\lambda_n(y)}-e^{it\lambda_n(x)}\|
                  \le\|\lambda_n(y)-\lambda_n(x)\|.
\]
The small link logarithm and fixed color-norm comparisons consequently
give
\begin{equation}
 \sup_{0\le t\le1}
   \left\|\frac1i\log\bigl(U^{e^{it\lambda_n}}\bigr)\right\|_\infty
                    \le C_g\rho.
 \label{csgq-path-margin}
\end{equation}
Here $C_g$ is independent of support, truncation radius, refinement and
volume. Condition \eqref{csgq-margin-condition} places these paths
strictly below $b/2$. At a smaller fixed endpoint ratio, restricting
the data radius gives a local germ instead; a whole-ball margin is
not inferred without the condition.

The input field in \eqref{csgq-path-margin} is summable rather than
compact. Before invoking the defining invariance of the quotient,
truncate its potential to a much larger cube containing the support of
$\lambda_n$ and all adjacent bonds. Inside the gauge support this
compact input is unchanged. Outside, the gauge is the identity and
the field is either unchanged or zero. The same strict path bound
therefore holds. Compact-field local gauge invariance and then
\eqref{csgq-summable-continuity}, as the input truncation tends to
infinity, give
\begin{equation}
              \Phi_E(\log U^{h_n}/i)=\Phi_E(\log U/i).
 \label{csgq-compact-path-identity}
\end{equation}
Now send $n$ to infinity using \eqref{csgq-compact-gauge-limit} and
\eqref{csgq-summable-continuity}. Extend constant global adjoint
invariance to summable fields by the same compact-input approximation.
The result is
\begin{equation}
 \Phi_E(a_L(w))
   =\Phi_E(\operatorname{Ad}_{c^{-1}}a_L(w^v))
   =\Phi_E(a_L(w^v)).
 \label{csgq-real-identity}
\end{equation}
This proof uses the same gauges and margins for every invariant class.

\subsection{Analytic continuation and precise scope}
\label{csgq-scope-section}

To finish Theorem \ref{csgq-main}, differentiate
\eqref{csgq-real-identity} at zero compact coarse gauge parameter.
The resulting infinitesimal identity is holomorphic in the finitely
many coarse field variables on their connected product ball and
vanishes on a real open set. It therefore vanishes throughout that
complex ball. Integrating it along any allowed compact local gauge
path that stays in the ball proves \eqref{csgq-gauge-identity}.
This argument does not extend an infinitesimal Ward identity to
arbitrary disconnected endpoint transformations.

The same endpoint method preserves the spatial symmetries required
by the invariant quotient, when the actual finite blocking and
minimum problem have those symmetries. The transformed ordinary
minimizer and the selected minimizer at transformed data are small
feasible fields with equal minimum values. The actual equality proof
supplies the controlled gauge; prefix covariance gives the endpoint
orbit relation; compact data give exponential tails. The preceding
compact-gauge argument applies, and the input total's spatial symmetry
supplies the remaining change of variables. Constant adjoint covariance
is exact for the equivariant construction, or follows from the small
comparisons and connectedness of the compact group. Symmetry of
individual carrier coefficients is unnecessary.

Universality means that every admissible interaction class can be
evaluated using this one fixed joint infinite background system.
It does not assert uniqueness among different regulator subsequences.
Semigroup identities require the maps and their compositions to
belong to the same joint system; finite-regulator identities alone
do not establish consistency for unrelated limits.

The connected representation and boundedness of the quotient action
remain separate analytic inputs. This section proves their missing
universal gauge identity from the actual finite transfer and flow,
endpoint decay, and summable-field continuity. It concerns homogeneous
small fields, $d\ge2$, and compact local gauge paths. It proves neither
quantum-integral descent on the quotient nor quantum induction,
large-orbit or mixed-history control, or a continuum measure.


% No preamble. Prefix pid-. Canonical normalized periodization only.
\section{Bounded finite periodization of analytic interaction classes}
\label{pid-section}

Use the infinite activity space and normalized-total quotient of
Notebook29. Fix the cell size $J$, at most $D$ scalar coordinates
per cell, exponent $a>0$, and radius $b>0$. Carriers are nonempty
finite connected subsets of the cell graph $\mathbb Z^d$, and
\begin{equation}
 \|E\|_{a,b}=\sup_i\sum_{H\ni i}e^{a|H|}\|E_H\|_b,
 \qquad
 \Phi_E(A)=\sum_H[E_H(A|_H)-E_H(0)]
 \label{pid-space}
\end{equation}
on compact fields. Let $\Lambda_N=(\mathbb Z/N\mathbb Z)^d$;
the underlying current lattice has period $NJ$. Write $\pi_N$
for cell and coordinate projection and $A^{\rm per}$ for the
infinite periodic lift of a finite field. It is bounded, generally
not summable. Internal link directions and color types are
identified by their actual periodic conventions.

\begin{theorem}[Canonical normalized periodization]
\label{pid-theorem}
Every class with a $J$-cell-translation-invariant normalized total
has a translation-covariant representative without increasing the
norm of any chosen input representative. There is a canonical
linear periodization map
\begin{equation}
 \operatorname{Per}_N:\mathcal Q^{\rm tr}_{a,b}
                   \longrightarrow\mathcal Q^N_{a,b},
 \qquad \|\operatorname{Per}_N\|\le1,
 \label{pid-operator}
\end{equation}
where the finite quotient uses connected torus carriers and
equality of vacuum-normalized totals. Its scalar functional is
holomorphic on the original open radius-$b$ polydisc. Every
fixed derivative tensor is the exact relative periodization of
the infinite total tensor. Total spatial and color symmetries
pass whenever they preserve the torus. Local gauge invariance
passes on the specified common logarithmic chart.
\end{theorem}

The proof retains raw constants in carrier norms and normalizes
only the total. It does not identify the result with an unrelated
actual finite-regulator quantum output.

\subsection{Translation averaging without norm loss}
\label{pid-averaging-section}

Let $\tau_t$ translate carriers, their coordinate slots and their
functions by a cell-lattice vector $t$. It is an isometry of the
activity space. If $[E]$ is translation invariant, then
$[\tau_tE]=[E]$. For centered F\o lner cubes $F_m$, put
\begin{equation}
 E^{(m)}=\frac1{|F_m|}\sum_{t\in F_m}\tau_tE.
 \label{pid-Cesaro}
\end{equation}
Its norm is at most $\|E\|_{a,b}$, and for every fixed $s$,
\begin{equation}
 \|\tau_sE^{(m)}-E^{(m)}\|_{a,b}
 \le\frac{|(F_m+s)\mathbin\triangle F_m|}{|F_m|}
                                      \|E\|_{a,b}\longrightarrow0.
 \label{pid-Folner-bound}
\end{equation}
Each fixed carrier has finitely many input coordinates. Countable
Montel extraction gives locally uniform limits $E_H^{\rm cov}$
for all carriers. Fatou, first on finite carrier collections and
smaller polydiscs and then at the original open radius, yields
\begin{equation}
 \|E^{\rm cov}\|_{a,b}\le\|E\|_{a,b}.
 \label{pid-covariant-bound}
\end{equation}
Equation~\eqref{pid-Folner-bound} passes componentwise and gives
exact covariance. On a compact input supported in $K$, the
omitted carriers with $|H|>Q$ contribute at most
$2|K|\|E\|_{a,b}e^{-aQ}$. The finite remaining sum converges,
so $\Phi_{E^{\rm cov}}=\Phi_E$. No uniqueness of this covariant
representative is needed.

Only cell translations are averaged. Finer unit translations,
point symmetries and color invariance need not hold for the
individual representative and will pass through its total.

\subsection{The exact orbit count and wrapped carrier norm}
\label{pid-orbit-section}

For a covariant representative define
\begin{equation}
 U_N^E(A)=\sum_{i\in\Lambda_N}\sum_{H\ni i}
                 \frac{E_H(A^{\rm per}|_H)}{|H|},
 \qquad
 \Phi_N^{[E]}(A)=U_N^E(A)-U_N^E(0).
 \label{pid-density}
\end{equation}
Choose one fundamental representative for each cell $i$. The
formula is independent of that choice by covariance. Its series
is absolutely convergent in the supremum norm on the full open
$b$ polydisc, bounded by $N^d\|E\|_{a,b}$ before normalization.

Let $\mathcal O$ be an orbit of finite carriers under
$N\mathbb Z^d$. A nonempty finite carrier has no nonzero
translation stabilizer. Every point in a representative $H$
corresponds to exactly one incidence pair $(i,H')$ with
$i\in\Lambda_N$, $H'\in\mathcal O$ and $i\in H'$. Hence
\begin{equation}
 \sum_{i\in\Lambda_N}\sum_{\substack{H'\in\mathcal O\\ i\in H'}}
                      \frac1{|H'|}=1.
 \label{pid-incidence}
\end{equation}
Covariance and periodicity make the summand value constant on
the orbit. Therefore
\begin{equation}
 U_N^E(A)=\sum_{\mathcal O}E_H(A^{\rm per}|_H),
                   \qquad H\text{ any representative of }\mathcal O.
 \label{pid-orbit-sum}
\end{equation}
This counts each carrier orbit once, including those whose
projection is not injective.

For a connected nonempty torus carrier $K$, set
\begin{equation}
 E_{N,K}(A|_K)=
       \sum_{\mathcal O:\,\pi_NH=K}E_H(A^{\rm per}|_H).
 \label{pid-grouped-family}
\end{equation}
Coordinate repetition is a linear map of maximum norm one.
Every summand is holomorphic on the original $b$ polydisc and
strictly supported on $K$. Projection preserves connectedness.

Fix a torus cell $j$ and one infinite lift. For every orbit
meeting $j$, choose one representative containing that lift.
Distinct orbits give distinct carriers. Since
$|\pi_NH|\le|H|$, this gives
\begin{align}
 \sum_{K\ni j}e^{a|K|}\|E_{N,K}\|_b
 &\le\sum_{\mathcal O:\,j\in\pi_NH}
                         e^{a|\pi_NH|}\|E_H\|_b\notag\\
 &\le\sum_{H\ni j}e^{a|H|}\|E_H\|_b
 \le\|E\|_{a,b}.
 \label{pid-wrapped-bound}
\end{align}
This proves normal convergence of the grouped family and the
no-loss carrier estimate. Choosing one incidence here is an
upper bound for the norm; the exact energy multiplicities were
already counted in \eqref{pid-incidence}. No factor
$|\pi_NH|/|H|$ belongs in the energy.

The raw grouped coefficients retain their constants. Normalizing
each coefficient separately could cost a factor two, and is
unnecessary. For example the constant singleton family
$E_{\{i\}}=c$ is the zero normalized class, but its raw periodic
total is $N^dc$. A canonical unnormalized exponential therefore
needs an additional vacuum convention.

\subsection{Derivatives, representative independence and local stability}
\label{pid-kernel-section}

Let $K_n^E(i_1,\alpha_1;\ldots;i_n,\alpha_n)$ be the infinite
total derivative kernel, with cell indices and internal types
separated. At every fixed $n$, connected-carrier Cauchy estimates
give absolute anchored summability with a positive spatial
exponent. Constants can depend on $n,D,J,a,b$; no uniform
all-order absolute Taylor norm is being asserted.

Differentiate \eqref{pid-orbit-sum}. A derivative in a periodic
coordinate differentiates each old occurrence of that coordinate.
Translate its first occurrence $i_1+Nz_1$ to a fixed lift $i_1$.
Every carrier through $i_1$ occurs exactly once. The remaining
slots then range over all their periodic lifts. Thus
\begin{align}
 &D^nU_N^E(0)[e_{i_1,\alpha_1},\ldots,e_{i_n,\alpha_n}]
 \notag\\
 &\qquad=\sum_{z_2,\ldots,z_n\in\mathbb Z^d}
 K_n^E(i_1,\alpha_1;
       i_2+Nz_2,\alpha_2;\ldots;i_n+Nz_n,\alpha_n).
 \label{pid-kernel-periodization}
\end{align}
Repeated finite slots include both coincident and distinct
infinite lifts with their ordinary multilinear multiplicities.
Absolute anchored summability justifies the rearrangement.

A null class has all $K_n=0$ for $n\ge1$. Its normalized
periodic total has every Taylor coefficient zero and hence is
zero on the connected finite-dimensional $b$ polydisc. Two
covariant representatives of the same class therefore give the
same normalized energy. This proves independence of the
translation-averaging subsequence. The scalar map is linear
even if representative selection is not. Taking representatives
approaching the quotient infimum in
\eqref{pid-covariant-bound},\eqref{pid-wrapped-bound} proves
\eqref{pid-operator}, and
\begin{equation}
 |\Phi_N^{[E]}(A)|\le2N^d\|[E]\|_{\mathcal Q}.
 \label{pid-energy-bound}
\end{equation}

There is exact scalar tiling consistency. For an integer $q>0$,
repeat $A$ on $\Lambda_{qN}$. The density in
\eqref{pid-density} is then $N$-periodic, so
\begin{equation}
 \Phi_{qN}^{[E]}(A\text{ repeated})=q^d\Phi_N^{[E]}(A).
 \label{pid-tiling}
\end{equation}
A rectangular multiple has its corresponding volume factor.

There is also eventual exact stabilization of each fixed lifted
coefficient. Let $K$ be a fixed finite connected infinite carrier
and choose, conservatively, $N>4\operatorname{diam}(K)+4$.
The copies $K+Nz$ are disjoint and nonadjacent. A connected $H$
with $\pi_NH=\pi_NK$ lies in their union, hence in one copy;
surjectivity of its projection makes it equal to that copy.
Equation~\eqref{pid-grouped-family} consequently has one orbit,
and under lifted coordinates
\begin{equation}
 E_{N,\pi_NK}=E_K
 \label{pid-fixed-carrier}
\end{equation}
on the full original polydisc. This gives local normal convergence
of canonical periodizations without a winding-tail estimate for
fixed $K$. It asserts neither global norm convergence nor
stabilization on carriers growing with $N$.

\subsection{Local gradients and exact finite gauge invariance}
\label{pid-gauge-section}

For a fixed infinite scalar coordinate $i$ and any bounded field
$\|A\|_\infty\le\rho<b$, define
\begin{equation}
 G_i^E(A)=\sum_{H\ni\operatorname{cell}(i)}
                                 \partial_iE_H(A|_H).
 \label{pid-gradient}
\end{equation}
Directional Cauchy gives
\begin{equation}
 |G_i^E(A)|\le\frac{\|E\|_{a,b}}{b-\rho}.
 \label{pid-gradient-bound}
\end{equation}
If $A,B$ agree in a $Q$-cell neighborhood of $i$, each differing
summand has a connected carrier of size at least $cQ$. Thus
\begin{equation}
 |G_i^E(A)-G_i^E(B)|
       \le\frac{2\|E\|_{a,b}}{b-\rho}e^{-caQ}.
 \label{pid-gradient-tail}
\end{equation}
This proves continuity under compact truncation of bounded
fields. On compact fields $G_i^E$ is the derivative of
$\Phi_E$. A null family therefore has $G_i^E=0$ on all
bounded interior fields, even without covariance of that
representative.

For a covariant representative, the same first-slot reanchoring
used above works at any periodic interior field and gives
\begin{equation}
 \partial_iU_N^E(A)=G_i^E(A^{\rm per}).
 \label{pid-periodic-gradient}
\end{equation}
This is another proof that null classes periodize to constant
raw totals and to zero normalized totals.

Suppose the infinite normalized total is locally gauge invariant.
For a fixed unit-current site $v$, its infinitesimal logarithmic
gauge field $\mathcal L_v(A)$ reads only the incident bond
potentials and is supported on those finitely many bonds.
On compact fields the Ward identity is
\begin{equation}
 \sum_{i\text{ incident to }v}
                    G_i^E(A)\mathcal L_{v,i}(A)=0.
 \label{pid-local-Ward}
\end{equation}
Compactly truncate an arbitrary bounded interior $A$. The local
generator values stabilize, and \eqref{pid-gradient-tail}
passes the gradient terms. Thus \eqref{pid-local-Ward} holds
for bounded periodic fields too. Using
\eqref{pid-periodic-gradient}, identify all incident orientations
and periodic multiplicities. The finite total obeys
\begin{equation}
 D\Phi_N^{[E]}(A)[\mathcal L_{\lambda,N}(A)]=0
 \label{pid-finite-Ward}
\end{equation}
for every finite torus gauge parameter. No noncompact infinite
gauge path has been applied to $\Phi_E$; only a local gradient
identity was extended by absolute convergence.

The logarithmic generator must be holomorphic on the domain.
If $b$ lies inside its fixed chart radius, the assertion holds
on the whole open $b$ polydisc. Otherwise keep full-$b$
holomorphy of the interaction and restrict the gauge assertion
explicitly to $b_g=\min(b/2,r_{\log}/2)$, where $r_{\log}$ is
a fixed radius of holomorphy for that generator. The actual
small-field setup already chooses $b$ inside this chart. Finite
gauge invariance follows along any gauge path whose fields stay
in the common domain. Large paths leaving it are not covered.

All finite nonlinear Ward equations also pass coefficientwise
in field degree. At any fixed order, differentiating the compact
infinite identity uses finitely many local generator coefficients,
nearest-neighbor shifts and coincident-slot embeddings of total
kernels. Equation~\eqref{pid-kernel-periodization} commutes with
these operations: summing images turns an infinite coincidence
condition into the exact periodic one. Fixed-order summability
justifies it. Equivalently, differentiate the exact analytic
identity \eqref{pid-finite-Ward}.

\subsection{Other symmetries and the finite-history boundary}
\label{pid-scope-section}

Any finer translation, hypercubic or global color symmetry of
the infinite total is inherited by its kernels. Relative
periodization commutes with each symmetry preserving the torus.
Taylor uniqueness gives the same symmetry of the normalized
finite total. Negative bonds use the signed reversal of their
positive-bond logarithmic coordinate. No symmetry of individual
$E_{N,K}$ is needed. Reality follows similarly; a real covariant
representative can also be obtained by averaging with its
complex-conjugate reflection without norm increase.

This completes Theorem~\ref{pid-theorem}. The canonical finite
functionals have exact Ward identities and a no-loss connected
representation. They need not equal an unrelated actual
finite-regulator family. Two uniformly bounded finite sequences
can have the same infinite local interaction while differing
by winding-supported terms at every finite size. Compactness
does not assert those terms are exactly absent.

Consequently, inserting $\operatorname{Per}_N[E]$ into a quantum
integral requires either an actual history/normalization comparison
or a construction using these periodized inputs from the outset.
Stationary backgrounds, retained windows, Gaussian normalizations
and gauge factors are not identified by periodizing the interaction
alone. Raw exponentials also require a separately specified vacuum
convention. No quantum integration identity, continuity of the
actual quantum update, or infinite induction is proved here.


% No preamble. Prefix qll-. Actual homogeneous local sector, fixed classical system.
\section{A local limit and quotient descent of the quantum macrostep}
\label{qll-section}

Fix the composite ratio \(R\), auxiliary grouping \(J\), preceding
homogeneous depth \(M\), real retained link/color windows, cutoff
threshold \(T\), positive coupling \(g\), and common field radius
\(b\) from the common-domain construction. All fluctuation-cell
dimensions are fixed once \(R,J\) are fixed. Compatible cubic
tori exhaust the current and next lattices together. Write
\[
 \|E\|_{a_0,b}
   =\sup_i\sum_{H\ni i}e^{a_0|H|}\|E_H\|_b,
 \qquad a_0=1/16,
\]
for the connected activity norm. Quotient classes identify
vacuum-normalized totals on compact fields, as in Notebook29.

The actual finite local correction has the decomposition
\begin{align}
 \mathcal F_N(E_N)
   &=\mathcal T_NE_N+\Delta_N+\mathcal Q_{{\rm det},N}
                         +\mathcal L_N(Q_{E_N}),\notag\\
 \mathcal G_N(E_N)
   &=\Delta_N+\mathcal Q_{{\rm det},N}
                         +\mathcal L_N(Q_{E_N}).
 \label{qll-finite-map}
\end{align}
Here \(\mathcal T_N\) is pre-gauge background substitution,
\(Q_{E_N}\) is the inherited fluctuation difference in the
actual chart, and \(\mathcal L_N\) is the strict Gaussian
cluster logarithm with its fixed classical base factors and
precision. The pivot term, determinant, sharp windows and
field-independent normalization are those of the finite
common-domain integral. Only after constructing this expression
do we subtract its value at zero background to form a quotient
total. A separately tracked vacuum normalization is not
identified with zero.

\subsection{One fixed, input-independent classical system}
\label{qll-data-section}

Call coefficients \emph{locally normal} when every fixed lifted
finite carrier coefficient converges uniformly on compact
subsets of its finite-dimensional background polydisc.
For joint coefficients require this convergence uniformly on
the fixed real window of their finitely many fluctuation
coordinates. Matrix coefficients use their finite block norm.
The final norm bound is on the original open \(b\) ball;
local normal convergence is only asserted on compact subsets.

Choose one common regulator subsequence with the following
actual classical data.

\begin{enumerate}
\item Every fixed baseline covariance entry \(C_{0,N}(i,j)\)
converges. The baseline has the uniform exponential Schur
bounds and positive spectral lower bound used in the strict
Gaussian theorem.
\item The strict precision coefficients \(B_{N,Y}\), classical
base exponents, pivot-log factors, and coherent pre-gauge
and joint-chart support coefficients converge locally normally.
Both matrix indices of \(B_{N,Y}\) belong to \(Y\).
\item All stronger support bounds, analytic neighborhoods,
strict sufficient budgets and exact local scalar identities
are those of the established finite construction. Aggregate
primitive labels by carrier and type before taking limits:
there are only finitely many types per fixed carrier, including
precision, base, inherited and bad-block types.
\end{enumerate}

These data are chosen independently of the inherited interaction.
At fixed parameters their strict families admit a countable
diagonal extraction. Each fixed coefficient has finitely many
coordinates. For the real fluctuation window, use the actual
unwindowed chart's holomorphic extension to a common complex
neighborhood and its uniform Cauchy bounds. Bounded measurable
functions alone would not give this compactness. The sharp
indicators remain fixed; they are not differentiated or subjected
to a complex change of variables.

The actual flat formulas already give compatible periodized
baseline kernels. More generally, joint entrywise extraction
with the stated uniform bounds preserves every finite
principal spectral inequality. The conclusions below refer to
the selected limiting classical system, not to an arbitrary
different sequence of baselines or classical branches.

\subsection{Explicit strict transport margins}
\label{qll-margins-section}

The common-domain proof permits the following stricter choices.
Fix its \(B=53/8\), primitive constants and classical support
reserves first. Choose \(R\) so that
\begin{equation}
 c_1(R)B\le a_0/4,
 \qquad c_1(R)=\frac{5^d2^d}{k_R},
 \qquad k_R=1+\lfloor(R-3)/4\rfloor .
 \label{qll-R-margin}
\end{equation}
After fixing \(R\), choose \(J\) and its logarithmic tail
collar so that, with the actual tail-to-margin ratio \(t\),
\begin{equation}
 t<1,\qquad
 D(J)\log\frac1{1-t}\le a_0/4.
 \label{qll-J-margin}
\end{equation}
This is possible because its sufficient bound is a constant
times \(J^{2d}e^{-\kappa_0\ell_{\rm tail}/2}\), which tends
to zero with the established choice of collar. It is independent
of the later choice of \(b\). Thus
\begin{equation}
 \delta_{\rm card}
   :=a_0-c_1(R)B-D(J)\log\frac1{1-t}
   \ge a_0/2>0.
 \label{qll-card-margin}
\end{equation}
The original weak half allocations alone would only give a
nonnegative margin; \eqref{qll-R-margin}--\eqref{qll-J-margin}
make the needed strictness explicit.

In the positive joint core/tail majorant, restricting original
carriers to \(|H|>Q\) now gives the bound
\begin{equation}
 C\epsilon e^{-\delta_{\rm card}Q}.
 \label{qll-input-carrier-tail}
\end{equation}
Choose also a fixed \(\zeta>1\) sufficiently close to one that
\begin{equation}
 \zeta t<1,\qquad
 c_1(R)B+D(J)\log\frac1{1-\zeta t}<a_0.
 \label{qll-tail-degree-margin}
\end{equation}
Replacing \(t\) by \(\zeta t\) in the exact scalar-multiindex
majorant gives an \(O(\zeta^{-k})\) bound above tail-expansion
degree \(k\). If \(t=0\), there are no positive tail degrees.
The rooted spatial weight and the strict spatial budget
\(\kappa_{\rm tail}>B+\delta\) control omitted remote supports.
The whole output-fiber maximum remains inside the support sum.

Choose \(b\), then the permitted \(T,g\), within the common-domain
order with strict slack in the precision, density, bad-block
and multiplier inequalities. Equivalently use the smaller
quarter budgets of the two-input comparison. Strictness allows
a fixed bookkeeping dilation greater than one in the positive
majorants below. All these choices are uniform in volume and
in preceding \(M\) when the established classical estimates
are uniform there.

\subsection{Continuity of the actual strict Gaussian construction}
\label{qll-Gaussian-section}

\begin{lemma}[Local construction continuity]
\label{qll-Gaussian-lemma}
Use fixed finite-dimensional fluctuation cells, fixed real
windows, the classical data above, and locally normal multiplier
coefficients under uniform strict budgets. Fix deterministic
connecting graphs that agree on every bounded lifted
configuration for sufficiently large periods. The canonical
connected Gaussian logarithm coefficients then have unique
locally normal limits for those limiting local data and
assignments. Their original anchored bound holds on the full
open \(b\) ball.
\end{lemma}

\begin{proof}
The proof retains the construction's positive majorants, not
merely its final activity norm.

First expand the precision inverse in its \(B_Y\) labels.
A term has the form
\[
 C_0(b,u_1)B_{Y_1}(u_1,v_1)C_0(v_1,u_2)
       \cdots B_{Y_r}(u_r,v_r)C_0(v_r,c).
\]
Every intermediate index belongs to a retained precision
carrier, because both indices of \(B_Y\) lie in \(Y\).
The endpoints \(b,c\) belong to the finite moment-root set
\(\mathcal M\). If \(X=\bigcup_jY_j\), the final connected
carrier contains \(\mathcal M\cup X\). Masked baseline
inverses act on the finite set \(\mathcal M\). External
determinant trace chains have the same property. Consequently
a fixed final carrier bounds every primitive spatial label
and covariance-chain index. No hidden unbounded index excursion
or new spatial reserve is needed. The full \(C_0\) kernel entry
is used as an entry, not as an unrecorded walk expansion.

Next keep tags on primitive positive weights and dilate them:
\begin{equation}
 B_Y\mapsto zB_Y,\qquad
 w_Y\mapsto zw_Y,\qquad q_i\mapsto zq_i,
 \qquad |z|\le1+\delta_0
 \label{qll-mark-dilation}
\end{equation}
for a sufficiently small fixed \(\delta_0>0\).
Here \(w_Y\) are interaction multipliers and \(q_i\) are
bad-block factors. Strict budgets make the same positive
overlap, endpoint-tree and Mayer majorants summable on this
disc. Scaling a bad-block factor only scales its joint-tail
activity weight; no new cutoff or differentiability assertion
is involved.

Every nonempty logarithm polymer contains an interaction or
bad-block mark. Mayer repetitions repeat such marks. Covariance
forest edges number at most the molecule count minus one,
while each inverse or determinant density order counts its
precision occurrences. Thus the remaining combinatorial tails
are \(O((1+\delta_0)^{-k})\), together with the strict geometric
precision-order tail if counted separately. There is no
unmarked unbounded combinatorial order left.

For a fixed final carrier and bounded orders there are only
finitely many discrete choices, since primitive labels have
already been aggregated by carrier and type. A retained term
depends on finitely many baseline entries and analytic
precision or multiplier coefficients. Its Gaussian moment
uses a finite principal covariance. The partition masks and
degree-adapted endpoint perturbations remain in the common
accretive neighborhood. Matrix determinants and inverses are
continuous holomorphic functions there. The established
density majorant is integrable on the entire fixed real
Gaussian space, uniformly in the regulator.

Dominated convergence therefore passes the Gaussian integrals
and compact forest-parameter integrals. The fixed sharp
indicators are only bounded measurable factors; they remain
untouched. Multiplier convergence is uniform on their compact
retained windows. The argument is uniform on compact background
polydiscs, or equivalently follows there from the common
holomorphic bounds and uniqueness of the limit.

The order-tail bounds now pass each complete coefficient.
Deterministic hull grouping preserves this convergence.
Winding configurations cannot remain in a fixed bounded
lifted carrier as the period grows. The final anchored
cardinality bound controls the omitted large carriers when
passing normalized totals on compact fields.

Finally apply Fatou to finite carrier selections at radius
\(r<b\), then increase \(r\) to \(b\) and exhaust the carriers.
This retains the original full open-domain bound. Compact-total
convergence follows by the same cardinality-tail argument as
analytic interaction compactness. The result is uniqueness of
the local limit for the selected local data and assignments,
not convergence in the global anchored norm.
\end{proof}

\subsection{Inherited insertions and external factors}
\label{qll-inherited-section}

Let finite input families converge locally normally with a
uniform activity budget. The core/tail construction evaluates
each input coefficient at an analytic core in a fixed strict
interior of its old \(b\) polydisc, then expands only in its
small supported tail. It retains the projected input carrier
and every attached background or fluctuation support.

Truncate original carrier size, remote support extent and
tail-expansion degree. Equations
\eqref{qll-input-carrier-tail}--\eqref{qll-tail-degree-margin}
and the rooted exponential bound give uniform positive tails.
On the remaining finite list, ordinary Cauchy estimates pass
all needed input derivatives and evaluations. The actual core
and fluctuation chart remain in their fixed interior margin,
uniformly on the real retained windows. The complete-fiber
norm prevents an extra output-coordinate count. Hence
\(Q_{E_N}\) has locally normal coefficient limits. The
compatible assignment is linear in the inherited input.
Equal-carrier exponentiation then preserves local convergence
and the required multiplier bounds.

Lemma~\ref{qll-Gaussian-lemma} applies to \(\mathcal L_N\).
The external pivot factors pass by local analytic composition.
The external determinant term passes by its absolutely
convergent closed-chain trace series, with every precision
label and its indices retained. Classical substitution passes
by the proved summable-field composition limit. We conclude
that \eqref{qll-finite-map} has locally normal connected
coefficient limits under these hypotheses.

\subsection{Canonical periodization and quantum quotient descent}
\label{qll-descent-section}

The periodization theorem supplies a covariant representative
of every translation-invariant class without norm increase,
a connected finite-periodic representation at the same
\(a_0,b\), and exact independence of that representative as
a finite vacuum-normalized total. It also supplies finite
local gauge and spatial identities on the common chart.
For a fixed covariant representative, its fixed lifted
coefficients stabilize exactly by
\eqref{pid-fixed-carrier}; in particular they converge locally
normally. On covariant representatives the orbit construction
is linear, including on differences.

\begin{theorem}[Quantum map on the selected local quotient]
\label{qll-descent-theorem}
Under the fixed classical-system, strict-margin and domain
hypotheses above, canonical periodized inputs define a
well-defined limiting map \(\mathcal F_g\) on a sufficiently
small ball of invariant analytic interaction classes.
Its output has a connected representation at the same
\(a_0,b\). The map preserves the local total gauge and spatial
identities of the actual sector. Writing \(\mathcal T\) for
classical substitution, the quantum remainder is
\(\mathcal G_g=\mathcal F_g-\mathcal T\).
\end{theorem}

\begin{proof}
Choose an allowed covariant representative \(E\). Its finite
periodizations and the fixed classical data meet the preceding
local-continuity hypotheses. They therefore produce a limiting
connected family \(\mathcal F(E)\).

If \(E'\) represents the same class, its finite periodized
normalized total is exactly the same. Their actual inherited
fluctuation differences agree as scalar insertions, and their
background terms agree after the chosen normalization.
If raw conventions differ, their constant scalar factor is
handled in the separate vacuum normalization. The fixed
Gaussian, windows and chart are otherwise identical.
Consequently their actual finite output normalized totals
agree, with compatible nonvanishing logarithm branches.
The canonical coefficient arrays need not agree.

Passing their local limits and then compact normalized totals
shows \([\mathcal F(E)]=[\mathcal F(E')]\). This proves descent.
The exact finite local gauge and spatial identities pass on
compact backgrounds and their allowed transforms. Holomorphic
continuation gives their local complex identities on the
common ball. Norm bounds pass by Fatou. The same reasoning
applies to \(\mathcal T\) and hence to the indicated remainder.
\end{proof}

\subsection{A quotient Lipschitz estimate without a high-jet hypothesis}
\label{qll-Lipschitz-section}

Let \(\bar\epsilon\) be a representative budget for which the
strict finite two-input comparison holds. It can be chosen so
the base exponent norm \(v\), inherited norms and bad-block
budget have the spare allocations used there; for example
\begin{equation}
 v\le\eta_*/(64e),\qquad
 q(E)\le\eta_*/(128e)
       \quad\hbox{when }\|E\|_{a_0,b}\le\bar\epsilon,
 \qquad 0<\eta_*\le1.
 \label{qll-comparison-budgets}
\end{equation}
The actual linear inherited difference satisfies
\begin{equation}
 q(D)\le16C_TC_R\,\frac{gT}{b}\|D\|_{a_0,b}.
 \label{qll-increment-bound}
\end{equation}
For two quotient classes of norm at most \(\bar\epsilon/4\),
choose covariant near-minimal representatives \(Y\) of \(y\)
and \(D\) of \(x-y\), and set \(X=Y+D\). For small \(\varepsilon>0\),
\begin{align}
 \|Y\|&\le\bar\epsilon/4+\varepsilon,\notag\\
 \|D\|&\le\|x-y\|_{\mathcal Q}+\varepsilon
                    \le\bar\epsilon/2+\varepsilon,\notag\\
 \|X\|&\le3\bar\epsilon/4+2\varepsilon<\bar\epsilon .
 \label{qll-compatible-representatives}
\end{align}
Covariant averaging permits these choices without norm loss.
This construction controls the representative difference;
two unrelated near-minimal representatives would not do so.

For completeness, interpolate the inherited exponent from
\(Q_Y\) to \(Q_Y+zQ_D\). When \(q(D)>0\), use
\begin{equation}
 L_D=\frac{\eta_*}{32e\,q(D)}\ge2.
 \label{qll-interpolation-radius}
\end{equation}
The strict base, inherited and bad-block budgets keep the
Gaussian logarithm family bounded by \(S_1\eta_*\) throughout
this disc. Banach-valued Cauchy gives the finite difference
bound \(64eS_1q(D)\). If \(q(D)=0\), the difference is zero.
Combining with \eqref{qll-increment-bound} yields, after
local limits, Fatou and \(\varepsilon\downarrow0\),
\begin{equation}
 \|\mathcal G_g(x)-\mathcal G_g(y)\|_{\mathcal Q}
 \le C_{\rm quant}\frac{gT}{b}\|x-y\|_{\mathcal Q},
 \qquad C_{\rm quant}=1024eS_1C_TC_R .
 \label{qll-quantum-Lipschitz}
\end{equation}
The pivot and determinant terms cancel because the classical
system is input independent. No fifth-order or other jet
vanishing of \(D\) was used in this logarithm comparison.

If \(B_R\) denotes the established bounded classical transport
constant at this fixed geometry, then
\begin{equation}
 \|\mathcal F_g(x)-\mathcal F_g(y)\|_{\mathcal Q}
 \le\left(B_R+C_{\rm quant}\frac{gT}{b}\right)
                                  \|x-y\|_{\mathcal Q}.
 \label{qll-full-Lipschitz}
\end{equation}
The finite comparison uses the same compatible carrier
assignments on both inputs. Its difference family converges
locally by the construction just proved, so Fatou applies
to that family before the quotient infimum is taken.
This is a small inherited-input dependence of the quantum
remainder; it is not contraction of the whole unextracted map.

\subsection{Varying actual inputs and finite-history compatibility}
\label{qll-varying-section}

\begin{theorem}[Varying-input compatibility]
\label{qll-varying-theorem}
Let finite input families \(E_N\) have exactly invariant
normalized totals, uniform strict budgets and a locally normal
limit \(E\). Use the same input-independent classical system,
fixed windows and exact insertion domains. The actual finite
outputs have local limits by the construction above, and their
limiting normalized class is \(\mathcal F_g([E])\).
The input and output carrier assignments need not be covariant.
\end{theorem}

\begin{proof}
Sections~\ref{qll-Gaussian-section} and
\ref{qll-inherited-section} apply directly to \(E_N\):
they require local input convergence and uniform positive
tails, not the designation \emph{canonical periodization}.
This produces a local output from that particular limiting
carrier family.

To identify its class, average each finite \(E_N\) over its
finite cell-translation group. The norm does not increase.
Its raw finite total is unchanged: normalized translation
invariance implies raw invariance because the zero field
is fixed by translation. Thus the actual finite quantum
output normalized total is also unchanged.

Extract a locally normal subsequence of these covariant
averages, with limit \(E^{\rm cov}\). Their identical finite
normalized totals and interaction compactness imply
\([E^{\rm cov}]=[E]\). The averaged inputs and the canonical
periodizations of \(E^{\rm cov}\) converge locally normally
to the same covariant family. Their input-independent
classical data are the same. Construction continuity therefore
gives identical canonical limiting output families.

Passing the equality between the original and averaged
finite output normalized totals identifies the original
limiting class with \(\mathcal F_g([E])\). This reasoning
applies to any extracted covariant average limit, so no
interaction-dependent replacement of the classical system
is needed. Canonical arrays themselves need not agree.
\end{proof}

By induction, every fixed finite number of actual homogeneous
local steps agrees, at the normalized local-total level, with
the corresponding composition of limiting quotient maps,
provided every intermediate actual finite input remains within
the uniform budgets and analytic insertion domains and the
classical systems and nested regulators are chosen jointly.
At each stage use Theorem~\ref{qll-varying-theorem}, then the
next stage's local continuity. Joint finite lists and countable
fixed parameter systems can be extracted by the earlier
diagonal construction when their hypotheses hold.

All intermediate windows, Gaussian integrations and generated
corrections remain present in this assertion. Each finite
vacuum normalization is tracked separately; the quotient
statement neither discards it from the integral nor proves
convergence of an extensive vacuum energy. This is not an
identification of sequentially cut off integrals with one
one-shot integral.

An infinite quotient relocalization does not by itself provide
small representatives of the actual finite winding remainders.
Thus an infinite quotient bound alone does not supply the
intermediate finite-input hypotheses just stated. If finite
classical reference subtraction is later used, the required
interface is local normal convergence of the actual reference
families and convergence of its extraction coefficients;
equality with a canonical periodized reference is unnecessary
and is not asserted here.

\subsection{Scope}
\label{qll-scope}

The result gives an actual local limit and a continuous
quantum map for the selected homogeneous local sector and
fixed classical system, with the established same-radius
analytic continuation and exact insertion representation.
It is not an assertion about the full Wilson measure or
mixed and large-field histories. It does not prove uniqueness
across arbitrary classical subsequences.

The one-input bounds and the small quantum Lipschitz factor
do not themselves provide an invariant renormalized ball.
A finite block still retains all generated quantum terms;
full-reference extraction, amplitude closure, an actual
coupling trajectory and the remaining field sectors and
observable limits require further work. No infinite
trajectory or continuum construction follows merely from
the finite-history compatibility proved here.


% No preamble. Prefix zif-. The optional fixed-coupling parity refinement stays in research.
\section{Quadratic forcing from a fixed outer field domain}
\label{zif-section}

Fix the composite geometry $R,J$, carrier exponent $a_0=1/16$,
and the jointly chosen, input-independent classical system of
the preceding local-limit theorem. The gauge Lie algebra is
semisimple, in particular it may be $\mathfrak{su}(2)$.
Constants may depend on the fixed geometry and color conventions,
but are uniform in preceding homogeneous depth $M$ and compatible
volume. The input of the quantum step in this section is zero.

\subsection{One outer construction and its restrictions}
\label{zif-outer-section}

The classical precision, stationary-map, pivot and base-interaction
coefficients have independent outer product domains. After $R,J$
are fixed, choose one radius $\rho>0$ on which all their strict
support and Gaussian-neighborhood budgets hold with a fixed
margin. No inherited activity is evaluated in this construction.
With the actual retained link/color window of threshold $T$, write
\begin{align}
 \mathcal G_{M,g,T}(0)
   &=\Delta_M+\mathcal Q_{{\rm det},M}
                         +\mathcal L_{M,g,T},\notag\\
 \mathcal L_{M,g,T}(w)
   &=\log\mathbb E_{C_M(w)}
             [\chi_T(\xi)e^{V_{\rm base}(w,\xi;g)}].
 \label{zif-zero-decomposition}
\end{align}
The covariance, stationary branch and charts are independent of
$g$. Field-independent powers of $g$, flat determinants, Haar
constants and the compact gauge-insertion normalization remain
in the separately tracked vacuum factor.

The actual outer-domain construction gives an anchored family
with
\begin{align}
 \|\mathcal G_{M,g,T}(0)\|_{a_0,\rho}
 &\le H(\rho,g,T),\notag\\
 H(\rho,g,T)
 &:=(C_\Delta+C_{\rm det})\rho
       +2S_1\bigl[e^3\tau(T)+e^{v(g,T)}v(g,T)\bigr],\notag\\
 v(g,T)&\le C_{\rm base}(gT^3+gT+g^2T^2),
 \label{zif-outer-bound}
\end{align}
with harmless fixed enlargements of the constants if needed.
Choose $T=T(g)\ge1$ and $g_0>0$ so that every strict base and
bad-block budget holds on $\rho$ for $0<g\le g_0$ and
\begin{equation}
 \sup_{M,\,0<g\le g_0}H(\rho,g,T(g))\le H_*<\infty.
 \label{zif-uniform-outer}
\end{equation}
A sufficiently large fixed multiple of
$\sqrt{\log(1/g)}$, with a fixed lower threshold, is one
permitted choice after reducing $g_0$. This is a bound on a
coefficient family, not on an extensive total.

For $0<b<\rho$, fix these outer data and restrict their
coefficients. The zero-input integral on $b$ is exactly the
restriction of \eqref{zif-zero-decomposition}: its chart,
covariance, window $T(g)$, measure factors and logarithm branch
are unchanged. A field radius enters the estimates, not the
defining integral. The infinite statement uses the same
outer-domain local-limit system followed by restriction.
Unrelated subsequences at different radii are not identified.
Two alternative valid carrier assignments for that same finite
integral also give the same normalized class by the established
local-limit theorem.

\subsection{Remove the total constant and linear sectors}
\label{zif-force-section}

\begin{theorem}[Uniform quadratic zero-input force]
\label{zif-theorem}
Under \eqref{zif-uniform-outer}, for every $b\le\rho/2$ the
zero-input correction has a representative of its
vacuum-normalized class satisfying
\begin{equation}
 \|[\mathcal G_{M,g,T(g)}(0)]\|_{\mathcal Q,a_0,b}
       \le C_F b^2,
 \qquad C_F=\frac{2H_*}{\rho^2}.
 \label{zif-force-bound}
\end{equation}
The constant is independent of $b,M$ and volume. The same
representative estimate holds before taking a local limit if
the hypotheses are imposed on the actual finite regulators.
\end{theorem}

\begin{proof}
The actual zero-input total is globally adjoint invariant:
the real retained-window change of variables, equivalently
the local gauge theorem of the common-domain integral, gives
this identity. Its derivative at zero in each link slot is
an invariant linear functional on the adjoint Lie algebra.
It vanishes on commutators, and
$[\mathfrak g,\mathfrak g]=\mathfrak g$; hence its total
linear jet is zero. This reasoning passes from finite
regulators to compact-field derivatives on the chosen limit.
Individual carrier coefficients need not have zero linear
jets or be separately gauge invariant.

Let $P_n$ be homogeneous phase projection in all background
coordinates of one carrier. Cauchy's formula on radii below
$\rho$, followed by their increasing limit, gives
\begin{equation}
 \|P_n E_H\|_b\le(b/\rho)^n\|E_H\|_\rho.
 \label{zif-phase-bound}
\end{equation}
Replace the outer family by
\begin{equation}
 E_H^{\ge2}:=\sum_{n\ge2}P_nE_H.
 \label{zif-representative}
\end{equation}
It retains every declared carrier and represents the same
normalized total: constants are normalized away and the
sum of the linear sectors is zero. Sum
\eqref{zif-phase-bound} first over degree for each carrier,
then over anchored carriers. This yields
\begin{equation}
 \|E^{\ge2}\|_{a_0,b}
 \le\frac{(b/\rho)^2}{1-b/\rho}\|E\|_{a_0,\rho}
 \le\frac{2H_*}{\rho^2}b^2.
 \label{zif-geometric-bound}
\end{equation}
There is no exchange of a root supremum with an order sum.
The argument is identical at finite volume and on the
normally convergent infinite family.
\end{proof}

For a Lie algebra with an abelian factor, the adjoint argument
alone does not remove a linear total; its vanishing would be
an additional hypothesis. Vacuum removal here changes only
the normalized class. It does not discard the corresponding
physical scalar normalization.

Let the actual full reference satisfy
$\beta(S_{MR})=1$ and
$\|[S_{MR}]\|_{\mathcal Q,b}\le C_Sb^2$, as obtained
from its fixed outer classical representation. Then
\begin{equation}
 P_{MR}=I-S_{MR}\beta,
 \qquad \|P_{MR}\|\le p:=1+C_\beta C_S
 \label{zif-projection}
\end{equation}
has a bound independent of $b$. Therefore
\begin{align}
 \|P_{MR}\mathcal G_{M,g,T(g)}(0)\|_{\mathcal Q,b}
      &\le C_0b^2,\qquad C_0=pC_F,\notag\\
 |\beta(\mathcal G_{M,g,T(g)}(0))|
      &\le C_\beta C_F.
 \label{zif-extracted-beta}
\end{align}
The $b^{-2}$ factor in the norm of $\beta$ is present;
the proven $b^2$ forcing cancels it. The finite representative
estimate does not by itself define a finite extraction
functional: a finite $\beta_N$ must separately have the
corresponding bounded normalization on the actual reference.

\subsection{Order of choices and the limits of the conclusion}
\label{zif-scope-section}

The outer radius and bound are chosen once after $R,J$.
The final inner radius may subsequently be reduced to meet
any fixed finite block of classical transport constraints.
For a prescribed family ceiling $\bar\epsilon$, the inherited
quantum difference coefficient is proportional to $gT(g)/b$.
A sufficient budget, with the constants of the common-domain
theorem, is
\begin{equation}
 \frac{16C_TC_RgT(g)}b\,\bar\epsilon
       \le\frac{\eta_*}{128e}.
 \label{zif-inherited-budget}
\end{equation}
The denominator $128e$ is the stronger two-input comparison
budget used for the quotient Lipschitz estimate. It is met by
reducing $g_0$ after the positive $b$ and
$\bar\epsilon$ have been fixed. Thus a ceiling of the form
$\bar\epsilon\le\eta_*/(\text{constant}\,a_Q)$,
where $a_Q=16C_TC_RgT/b$, creates no radius cycle.
The base and bad-block budgets remain separate requirements.

If all intermediate extracted amplitudes have norm at most
$C_{\rm int}b^2$, the bounded classical transport, the force
bound and the fixed-coupling quantum Lipschitz bound give
an unprojected one-step amplitude bounded by a constant
times $b^2$. Its Maxwell increment is consequently bounded
independently of $b$. This provides neither its sign nor
control of infinitely many accumulated increments.

Finally, $\Delta_M$ and $\mathcal Q_{{\rm det},M}$ have
no small factor $g$. Their removal of constant and linear
sectors yields $O(b^2)$, not $O(g^2)$. A bound of the latter
kind for the full force would require a separate proved
cancellation or subtraction. Nothing here is a contraction
in the joint coupling--activity variables, a large-field
estimate, or an identification of the local sector with the
full Wilson measure.


% No preamble. Prefix eqr-. Infinite local quotient; finite-history boundary explicit.
\section{Reference extraction and a contractive finite-block recurrence}
\label{eqr-section}

Work on the selected joint homogeneous classical system of
Sections~\ref{csgq-section} and \ref{qll-section}. Fix the primitive
odd ratio $\ell$, the elementary composite ratio $R=\ell^s$, the
auxiliary grouping $J$, and carrier exponent $a_0=1/16$.
At stage $j$ the preceding depth is $M_j=M_0R^j$. Each stage uses
the same cell grouping and analytic radius in its own current units.
All constants below may depend on this fixed geometry, but are uniform
in $M_j$. Use a simple gauge Lie algebra, in particular
$\mathfrak{su}(2)$, so the Maxwell functional is scalar. The finite
family of componentwise references gives the corresponding version
for a semisimple product.

The local quantum quotient map is an established input from
Theorem~\ref{qll-descent-theorem}, with its strict domain, periodization,
classical-system and logarithm-branch hypotheses. The discussion below
does not replace that theorem by an inference from compact-field
equality alone. The actual finite-regulator history remains a separate
issue stated at the end.

\subsection{The actual reference transports exactly}
\label{eqr-reference-section}

Let $U_M(D)$ be the selected ordinary homogeneous minimum and let
\begin{equation}
 S_M(D)=S_{\rm bare}(U_M(D)),\qquad S_M(1)=0,
 \qquad
 P_{M,R}(W)=\mathcal M_\ell^k U_{MR}(W),\quad M=\ell^k.
 \label{eqr-reference-definition}
\end{equation}
The trace convention and plaquette counting are fixed, and the coupling
is outside $S_M$. Classical background substitution is
\[
        (\mathcal T_{M,R}E)(W)=E(P_{M,R}(W)).
\]
It contains no fluctuation integration.

\begin{lemma}[Exact reference identity]
\label{eqr-reference-lemma}
On the specified common local analytic domains in dimension four,
\begin{equation}
                  \mathcal T_{M,R}S_M=S_{MR}.
 \label{eqr-reference-transport}
\end{equation}
This identity holds for the actual ordinary classical reference,
not merely for a reference with the same quadratic normalization.
It passes to the jointly chosen infinite normalized classes.
\end{lemma}
\begin{proof}
On a finite compatible regulator put
\[
 U_2=U_{MR}(W),\qquad D=\mathcal M_\ell^kU_2,\qquad U_1=U_M(D).
\]
Both fine configurations are fields on the same oriented graph.
Its fine spacing is $1/M$ in $D$ units and $1/(MR)$ in $W$ units.
Exact ordinary composition gives
\begin{equation}
 \mathcal M_\ell^kU_1=D=\mathcal M_\ell^kU_2,\qquad
 \mathcal M_\ell^{k+s}U_1=W=\mathcal M_\ell^{k+s}U_2.
 \label{eqr-two-constraints}
\end{equation}
The actual minimum comparison balls can be made common for fixed $R$.
Indeed a fine physical potential changes units by
$A_1^W(x)=R A_1(Rx)$, so
\begin{equation}
 \|A_1^W\|_{C^{1,\alpha}}
       \le CR^{2+\alpha}\|A_1\|_{C^{1,\alpha}}.
 \label{eqr-reference-rescaling}
\end{equation}
The prefix bound is
$\|\log D/i\|_\infty\le C_A\|\log W/i\|_\infty/R$.
Thus, for fixed $R$, a sufficiently small real $W$ ball places
both fields in both needed comparison classes, uniformly in
preceding refinement and volume.

The minimum at $D$ gives $S_{\rm bare}(U_1)\le S_{\rm bare}(U_2)$;
the minimum at $W$ gives the reverse inequality. In dimension four
the coupling-free bare action is the same plaquette sum under this
coordinate relabeling: its physical prefactor is
$\eta^{d-4}=1$. Hence the two actions agree. This is also the identity
$S_M(V^\star)=S_M(V_p)=S_{MR}$ in the actual common-domain
factorization.

Both sides are holomorphic on the specified connected common domain.
Uniqueness from the real germ continues the identity there.
For $R>C_A$, the angle estimate maps the common $b$ ball into itself;
the smaller real comparison ball does not cause a repeated analytic
radius loss. No minimum assertion at complex fields is used.

For the fixed joint infinite system, pass vacuum-normalized values on
compact $W$ using the connected interaction bounds and summable-field
continuity. These are the same hypotheses used for the classical
substitution limit. No convergence of an extensive infinite action
value is asserted. The finite references vanish at the trivial field,
so the normalized identity has no vacuum ambiguity.
\end{proof}

The exact flat midpoint symbol and its identification with the
classical Hessian give
\begin{equation}
 q_M(p)=|p|^2I-pp^T+O(|p|^4),\qquad \beta(S_M)=1.
 \label{eqr-reference-beta}
\end{equation}
This normalization was proved by the flat symbol calculation;
ellipticity alone would not imply it. The fixed-outer-radius
classical support estimate and its zero total linear jet also give
\begin{equation}
                    \|[S_M]\|_{\mathcal Q,b}\le C_Sb^2,
 \label{eqr-reference-size}
\end{equation}
uniformly in $M$ and in the subsequent inner radius $b$.
An arbitrary Maxwell-normalized local Wilson reference need not obey
\eqref{eqr-reference-transport} as a full nonlinear functional.

\subsection{The projection telescope}
\label{eqr-projection-section}

Write $\mathcal Q_j$ for the invariant quotient at stage $j$ and set
\begin{equation}
 S_j=[S_{M_j}],\qquad
 P_j=I-S_j\beta_j,\qquad
 \mathcal T_j=\mathcal T_{M_j,R},\qquad
 A_j=P_{j+1}\mathcal T_j.
 \label{eqr-projections}
\end{equation}
The notation $S_j\beta_j$ denotes a rank-one operator.
Because $\beta_j(S_j)=1$, $P_j$ projects onto $\ker\beta_j$.
The established coefficient and reference estimates give
\begin{equation}
 \|\beta_j\|\le C_\beta b^{-2},\qquad
 \|P_j\|\le p:=1+C_\beta C_S.
 \label{eqr-projection-norm}
\end{equation}
The $b^{-2}$ factor has been cancelled by
\eqref{eqr-reference-size}, not omitted.

\begin{lemma}[Cancellation of every intermediate extraction]
\label{eqr-telescope-lemma}
The actual references satisfy
\begin{equation}
 P_{j+1}\mathcal T_jP_j=P_{j+1}\mathcal T_j.
 \label{eqr-intertwining}
\end{equation}
Consequently, for every finite $q$,
\begin{align}
 A_{j+q-1}\cdots A_j
   &=P_{j+q}\mathcal T_{j+q-1}\cdots\mathcal T_j,
       \label{eqr-product}\\
 A_{j+q-1}\cdots A_jP_j
   &=P_{j+q}\mathcal T_{M_j,R^q}P_j.
       \label{eqr-direct-product}
\end{align}
The second identity uses the classical semigroup on the same joint
invariant system.
\end{lemma}
\begin{proof}
For $x\in\mathcal Q_j$, the difference between the two sides of
\eqref{eqr-intertwining} is
$-\beta_j(x)P_{j+1}\mathcal T_jS_j$.
By \eqref{eqr-reference-transport}, this is
$-\beta_j(x)P_{j+1}S_{j+1}=0$.
Repeated application removes the intermediate projections.
The actual scalar semigroup modulo gauge, together with the
universal summable-gauge identity of Section~\ref{csgq-section},
gives \eqref{eqr-direct-product}.
\end{proof}

The proof does not require
$\beta_{j+1}\mathcal T_j=\beta_j$.
Nor does it require equality of sequential and direct carrier arrays.
The common scalar class may be estimated using its separately
justified direct connected representation. If instead
$\mathcal T_jS_j=S_{j+1}+D_j$, the error in
\eqref{eqr-intertwining} is
$-P_{j+1}D_j\beta_j$, potentially costing
$b^{-2}\|D_j\|$. The actual identity makes that error exactly zero.

Use the supported low-kernel and high-jet transport estimates on the
same infinite system. Their explicit common-domain and reserve
hypotheses give
\begin{equation}
 \left\|A_{j+q-1}\cdots A_j\big|_{\ker\beta_j}\right\|
 \le\theta_q:=
 p^2\left[C_{\rm low}(J)R^{-q\alpha}
                     +6C_{\rm high}R^{-q}\right].
 \label{eqr-classical-block}
\end{equation}
The fixed grouping $J$ is chosen before $q$; its finite constant is
not assumed uniform as $J$ grows. Choose a finite $q$ so that
$\theta_q\le1/4$. Then choose $b$ to meet the direct $R^q$
high-jet disk and all classical margins. The constants in
\eqref{eqr-classical-block} do not deteriorate just because $b$
is reduced. Bounded one-step transport gives constants
\begin{equation}
 \|\mathcal T_j\|\le B_T,\qquad
                 \|A_j\|\le B,\qquad B\ge\max(1,pB_T),
 \label{eqr-one-step-linear}
\end{equation}
uniform in $j$. These may be large; no one-step classical
contraction is asserted.

\subsection{The local quantum remainder at fixed coupling}
\label{eqr-quantum-section}

Let $\mathcal F_j$ be the actual limiting local quotient map with
specified coupling $g_j$, retained threshold $T_{{\rm cut},j}$,
and the fixed classical system. Write
\begin{equation}
 \mathcal F_j(x)=\mathcal T_jx+\mathcal G_j(x),\qquad
 \Psi_j(x):=P_{j+1}\mathcal F_j(x)=A_jx+H_j(x),
 \qquad H_j=P_{j+1}\mathcal G_j.
 \label{eqr-quantum-split}
\end{equation}
The remainder includes the pivot and Gaussian determinant terms and
the full base Gaussian logarithm. Those terms are not discarded.
The output of $\Psi_j$ belongs to $\ker\beta_{j+1}$.

Let $\bar\epsilon$ be the representative ceiling from the strict
two-input comparison. The quotient maps are defined on the ball of
radius $\bar\epsilon/4$. Equation
\eqref{qll-quantum-Lipschitz} gives
\begin{equation}
 \|\mathcal G_j(x)-\mathcal G_j(y)\|_{\mathcal Q}
 \le L_{{\rm quant},j}\|x-y\|_{\mathcal Q},\qquad
 L_{{\rm quant},j}
       =C_{\rm quant}\frac{g_jT_{{\rm cut},j}}b.
 \label{eqr-quantum-lipschitz}
\end{equation}
It has no equal-low-jet hypothesis. One can see the domain margin
directly from \eqref{qll-compatible-representatives}: choose
covariant near-minimal representatives for $y$ and $x-y$ and add
them. Both complete inputs then have norm less than
$\bar\epsilon$. Independent near-minimal representatives for $x$
and $y$ would not control their difference.

Choose a uniform number $\nu$ with
\begin{equation}
 pL_{{\rm quant},j}\le\nu,\qquad
 \|H_j(x)-H_j(y)\|\le\nu\|x-y\|.
 \label{eqr-projected-lipschitz}
\end{equation}
All comparisons here hold the coupling, threshold, classical saddle
and cutoff data fixed at each step. The limiting map is well defined
by Theorem~\ref{qll-descent-theorem}; the finite comparison and its
passage to the quotient are provided by
Section~\ref{qll-Lipschitz-section}.

There are two useful forcing estimates. The original common-domain
bound gives
\begin{align}
 \|H_j(0)\|&\le f, \notag\\
 f&\le p\sup_j\left[
 (C_\Delta+C_{\rm det})b+
 2S_1\{e^3\tau(T_{{\rm cut},j})+e^{v_j}v_j\}\right],
 \notag\\
 v_j&\le C_{\rm base}
       (g_jT_{{\rm cut},j}^3+g_jT_{{\rm cut},j}
                              +g_j^2T_{{\rm cut},j}^2).
 \label{eqr-crude-force}
\end{align}
More strongly, fix the outer zero-input construction of
Section~\ref{zif-section} and then restrict to $b\le\rho/2$.
Equations \eqref{zif-force-bound} and
\eqref{zif-extracted-beta} yield
\begin{equation}
 \|\mathcal G_j(0)\|\le C_Fb^2,\qquad
 \|H_j(0)\|\le C_0b^2,\qquad
 C_F=\frac{2H_*}{\rho^2},\quad C_0=pC_F.
 \label{eqr-quadratic-force}
\end{equation}
These are the same actual maps restricted from a common outer
domain, not radius-dependent replacements. The force need not be
$O(g_j^2)$.

\subsection{A finite-block perturbation estimate}
\label{eqr-block-section}

\begin{lemma}[Duhamel bound with a contractive classical block]
\label{eqr-Duhamel-lemma}
Suppose the maps in \eqref{eqr-quantum-split} are evaluated on
their common admissible neighborhood,
\eqref{eqr-classical-block} and \eqref{eqr-one-step-linear} hold,
and \eqref{eqr-projected-lipschitz} holds. For two trajectories
using the same prescribed step couplings and other fixed data,
with their initial difference in $\ker\beta_j$, one has
\begin{align}
 \|\delta_q\|
 &\le\left[
 \theta_q+\nu\sum_{k=0}^{q-1}
                    B^{q-1-k}(B+\nu)^k\right]\|\delta_0\|
       \notag\\
 &=\left[\theta_q+(B+\nu)^q-B^q\right]\|\delta_0\|
       \notag\\
 &\le\left[\theta_q+q\nu(B+\nu)^{q-1}\right]\|\delta_0\|.
 \label{eqr-Duhamel-bound}
\end{align}
The sum formula includes $\nu=0$ by continuity.
If $\|H_i(0)\|\le f$, then any trajectory obeys
\begin{equation}
 \|x_{j+k}\|\le(B+\nu)^k\|x_j\|
                     +f\sum_{l=0}^{k-1}(B+\nu)^l
 \label{eqr-elementary-size}
\end{equation}
as long as the intermediate maps remain admissible.
\end{lemma}
\begin{proof}
The exact variation-of-constants identity is
\begin{align}
 \delta_q={}&A_{j+q-1}\cdots A_j\delta_0 \notag\\
 &+\sum_{k=0}^{q-1}A_{j+q-1}\cdots A_{j+k+1}
   \bigl[H_{j+k}(x_{j+k})-H_{j+k}(y_{j+k})\bigr].
 \label{eqr-Duhamel-identity}
\end{align}
An empty product is the identity. The elementary difference bound
is $\|\delta_k\|\le(B+\nu)^k\|\delta_0\|$.
Apply the improved classical product bound only to the first term
in \eqref{eqr-Duhamel-identity}; bound each shorter product by the
corresponding power of $B$. This proves the first line of
\eqref{eqr-Duhamel-bound}. The difference-of-powers identity gives
the second line, and bounding every sum term by
$(B+\nu)^{q-1}$ gives the third.

For a single step,
$\|\Psi_i(x)\|\le(B+\nu)\|x\|+f$.
Induction proves \eqref{eqr-elementary-size}.
No contraction of a shorter classical product was presumed.
\end{proof}

The leading classical term in \eqref{eqr-Duhamel-bound} has no
growing factor from intermediate projections. The remaining
perturbative sum has finite length $q$, fixed before the coupling
is made small.

\subsection{An invariant quadratic amplitude envelope}
\label{eqr-envelope-section}

\begin{theorem}[Extracted finite-block contraction on the local quotient]
\label{eqr-envelope-theorem}
Use the selected joint infinite homogeneous system, the actual
reference family, the transport estimates
\eqref{eqr-classical-block}--\eqref{eqr-one-step-linear},
the quantum maps of Theorem~\ref{qll-descent-theorem}, and the
fixed-outer-radius force of Theorem~\ref{zif-theorem}.
Choose a finite $q$ with $\theta_q\le1/4$. Put
\begin{equation}
 K=B+1,\qquad L_q=\sum_{l=0}^{q-1}K^l,\qquad
 K_{\rm amp}=\max\{1,2C_0L_q\},\qquad
 C_{\rm int}=K^qK_{\rm amp}+C_0L_q.
 \label{eqr-envelope-constants}
\end{equation}
Choose the common positive $b\le\rho/2$ to satisfy all direct
$R^q$ classical and high-jet domain restrictions and
\begin{equation}
                  C_{\rm int}b^2<\bar\epsilon/4.
 \label{eqr-inner-radius}
\end{equation}
Finally choose the weak-coupling range and thresholds so that the
strict two-input budgets, including
\eqref{zif-inherited-budget}, hold and
\begin{equation}
             \nu\le\min\left\{1,\frac1{4qK^{q-1}}\right\}.
 \label{eqr-nu-choice}
\end{equation}

For each initial stage $j$, every $q$-step extracted block maps
\begin{equation}
 \{x\in\ker\beta_j:\ \|x\|\le r_b\},
                \qquad r_b=K_{\rm amp}b^2,
 \label{eqr-invariant-ball}
\end{equation}
into the corresponding radius-$r_b$ ball at stage $j+q$, with
Lipschitz constant at most $1/2$. Every intervening amplitude is
at most $C_{\rm int}b^2<\bar\epsilon/4$.

For a prescribed countable sequence of couplings in this weak
range, with its fixed common classical local-limit system,
the extracted quotient recursion is therefore defined for all
stages starting in \eqref{eqr-invariant-ball}. At block endpoints
two trajectories using that same sequence contract by $2^{-n}$
after $n$ blocks. These conclusions also bound a density-preserving
running trajectory for as long as it stays in the prescribed
weak-coupling regime.
\end{theorem}
\begin{proof}
Start in \eqref{eqr-invariant-ball}. For $0\le k\le q$,
\eqref{eqr-elementary-size}, $\nu\le1$, and
\eqref{eqr-quadratic-force} give
\begin{equation}
 \|x_{j+k}\|
 \le K^kK_{\rm amp}b^2+C_0b^2\sum_{l=0}^{k-1}K^l
 \le C_{\rm int}b^2<\bar\epsilon/4.
 \label{eqr-intermediate-bound}
\end{equation}
This induction verifies, before the next step is invoked, that
each input lies in the common admissible neighborhood. The
same holds for the trajectory starting at zero and for any
second input in the ball.

Equation \eqref{eqr-Duhamel-bound} and
\eqref{eqr-nu-choice} make the block Lipschitz constant at most
$1/4+1/4=1/2$. Its zero-input output has norm at most
$C_0L_qb^2\le r_b/2$. Therefore its output at any $x$ in
\eqref{eqr-invariant-ball} has norm at most
$\|x\|/2+r_b/2\le r_b$.
The output is extracted at every step by definition.

All estimates are uniform in the initial homogeneous depth.
Repeat the argument block by block. This constructs the
stated infinite sequence of quotient classes for fixed
prescribed step data; it asserts no convergence of a single
nonautonomous trajectory to a fixed point. Applying the same
block Lipschitz bound to two trajectories with identical data
gives the $2^{-n}$ difference estimate.
\end{proof}

The parameter order is noncircular. Choose the primitive data,
the elementary $R$ and grouping $J$ with the strict
Section~\ref{qll-margins-section} margins, then choose the finite
$q$ using \eqref{eqr-classical-block}. The outer radius and
$H_*$ in Section~\ref{zif-section} are fixed before the final
inner radius. Then choose $b$ by
\eqref{eqr-inner-radius} and the direct $R^q$ restrictions.
Only afterward decrease the coupling ceiling. For example,
thresholds equal to a sufficiently large fixed multiple of
$\sqrt{\log(1/g)}$, with a fixed lower threshold, have
$gT(g)/b\to0$, $gT(g)^3\to0$ and $\tau(T(g))\to0$.
Thus the comparison and forcing budgets can be met uniformly
on $0<g\le g_0$. In particular
\[
 \frac{16C_TC_RgT(g)}b\,\bar\epsilon
                   \le\frac{\eta_*}{128e}
\]
is imposed with the stronger two-input budget.
The constants $p,B,\theta_q$ do not grow just because $b$
is reduced.

For comparison, the weaker force \eqref{eqr-crude-force} also
suffices for a fixed-radius envelope. With the same $K,L_q$,
take
\begin{equation}
 r=\frac{\bar\epsilon}{8K^q},\qquad
 f\le\frac{r}{2L_q}.
 \label{eqr-crude-envelope}
\end{equation}
Then intermediate norms are at most
$\bar\epsilon/8+r/2<\bar\epsilon/4$, and the same proof
preserves the radius-$r$ extracted ball. The quadratic
version is stronger because it controls the Maxwell
increment independently of the inner radius.

\subsection{Density bookkeeping and bounded coupling increments}
\label{eqr-coupling-section}

Along the quadratic envelope the unprojected one-step value obeys
\begin{equation}
 \|\mathcal F_j(x_j)\|
 \le \left[B_TC_{\rm int}+C_F
                       +L_{{\rm quant},j}C_{\rm int}\right]b^2.
 \label{eqr-unprojected-size}
\end{equation}
This follows by separating $\mathcal T_jx_j$,
$\mathcal G_j(0)$ and
$\mathcal G_j(x_j)-\mathcal G_j(0)$.
The norm of $\beta$ consequently gives
\begin{equation}
 |\beta_{j+1}(\mathcal F_j(x_j))|
 \le C_\beta\left[B_TC_{\rm int}+C_F
                       +L_{{\rm quant},j}C_{\rm int}\right].
 \label{eqr-beta-increment-bound}
\end{equation}
The factor $b^{-2}$ has again been paid by a proved $b^2$
amplitude. This is an absolute bound, with no sign.

For the actual density, write $\gamma_j=g_j^{-2}$. The
one-step classical factorization is
\[
        \exp[-\gamma_j S_{j+1}+\mathcal F_j(x_j)].
\]
Full-reference extraction gives the exact identity
\begin{equation}
 \mathcal F_j(x_j)
  =\beta_{j+1}(\mathcal F_j(x_j))S_{j+1}+\Psi_j(x_j),
 \qquad
 \gamma_{j+1}
  =\gamma_j-\beta_{j+1}(\mathcal F_j(x_j)).
 \label{eqr-density-running}
\end{equation}
The negative sign is forced by the displayed classical exponent.
A freely prescribed weak coupling sequence need not satisfy
\eqref{eqr-density-running}. The theorem bounds the extracted
amplitudes uniformly over such sequences as operator estimates;
it identifies a continuation of the same density only along
a sequence satisfying this extra update.

Two trajectories that separately update their couplings need
not use the same maps: their coupling and threshold parameters
also differ. Equation \eqref{eqr-quantum-lipschitz} controls only
the inherited-input dependence at fixed coupling. Thus the
$1/2$ block estimate is not a contraction in the joint
coupling--activity variables.

For a fixed finite number of steps,
\eqref{eqr-beta-increment-bound} permits choosing the initial
inverse coupling sufficiently large to retain a positive
weak-coupling range. It does not control an infinite sum of
increments, establish their sign, or prove asymptotic freedom.
It does not prove that a physical running trajectory remains
in the regime needed by Theorem~\ref{eqr-envelope-theorem}.

\subsection{What the quotient envelope does not supply}
\label{eqr-scope-section}

The conclusions concern the invariant infinite quotient of
the actual homogeneous local germ and its specified analytic
continuation. They hold for every admissible input class
on one fixed joint classical system. No uniqueness across
arbitrary classical regulator subsequences is asserted.

Theorem~\ref{qll-varying-theorem} identifies fixed finite
histories only when their actual finite intermediate input
families satisfy the uniform budgets and insertion-domain
hypotheses. A small quotient norm, or a newly localized
infinite representative, does not automatically provide those
bounds for the actual finite winding remainders. For a fixed
finite $q$, crude canonical estimates and an additional radius
reduction may furnish a finite history, but that construction
must be supplied separately. The infinite quotient envelope
is not a regulator-uniform finite RG induction.

If extraction is performed before the local limit, it must
use the actual finite reference $S_{M,N}$ and a specified
finite coefficient functional $\beta_N$ with the required
normalization and convergence. The finite-volume Ward
construction supplies a candidate cutoff-jet functional.
One must not assume that $S_{M,N}$ equals the canonical
periodization of an independently selected infinite reference.
The exact classical identity holds for the actual finite
references; compatibility of projected finite histories
requires the separate coefficient and family statements.

The quotient removes field-independent totals. Physical
vacuum factors remain separately tracked, and their
extensive limits are not supplied by this theorem.
Large-field and mixed histories, global measure consistency,
observables and the continuum limit remain outside this
local amplitude result.

The intermediate projection obstruction is therefore resolved:
the actual classical reference transports exactly, so its
extractions telescope. Combined with the proved local quantum
quotient map and quadratic force, this yields the finite-block
contraction and invariant amplitude envelope above. It does
not constitute a completed Yang--Mills construction.


\section{The next construction problem}
\label{qdn-next-section}

The normalized local quantum map and its finite-history compatibility
now have a specified construction. Their domains and positive
majorants are uniform in preceding homogeneous classical depth.
They provide the setting in which an equivalent connected
representation can be used at the next local integration.
Field-independent vacuum factors remain separately tracked;
no convergence of an extensive vacuum energy is inferred.

The actual inverse coupling obeys
\[
 \gamma_{j+1}=\gamma_j-
       \beta_{j+1}\bigl(\mathcal F_{j,g_j}(E_j)\bigr),
 \qquad \gamma_j=g_j^{-2}.
\]
Completing its weak-coupling trajectory is the next essential
step. Absolute increment bounds alone do not establish its
direction or permit the ultraviolet cutoff to be removed at a
fixed interacting endpoint. The generated drift must be evaluated
with the actual determinants, cutoff factors and inherited
interaction. Its normalization must agree with the reference
action used in the exact density identity.

The activity estimates concern the specified homogeneous
small-field sector. The full Wilson construction must also
incorporate mixed and large-field histories, identify the
continued local representation on their required gauge-orbit
domains, and obtain uniform estimates with observable sources.
Those estimates feed the compactness, rotation, nontriviality
and reconstruction obligations in Notebook 01.

\begin{thebibliography}{9}
\bibitem{priorqdn}
\emph{Toward a four-dimensional Yang--Mills construction},
research Notebooks 15--29 in this workspace. The ordinary minimum
and gauge transfer, strict overlapping Gaussian expansion,
common-domain quantum map and total Ward extraction are the
principal inputs.
\bibitem{balaban109qdn}
T.~Ba\l aban,
\emph{Renormalization group approach to lattice gauge field
theories. I. Generation of effective actions in a small field
approximation and a coupling constant renormalization in four
dimensions},
Commun. Math. Phys. \textbf{109} (1987), 249--301.
\href{https://doi.org/10.1007/BF01215223}{Publisher record}.
\bibitem{balaban116qdn}
T.~Ba\l aban,
\emph{Renormalization group approach to lattice gauge field
theories. II. Cluster expansions},
Commun. Math. Phys. \textbf{116} (1988), 1--22.
\href{https://doi.org/10.1007/BF01239022}{Publisher record}.
The finite localization and resummation architecture is checked
against the original article. The quotient and local-continuity
proofs in this notebook are stated and proved separately.
\bibitem{bkarqdn}
A.~Abdesselam, notes on the BKAR formula, Theorem 1.
\href{https://mabdesselam.github.io/BKAR.pdf}{Author-hosted notes}.
The finite forest identity is used with the explicit positive
majorants and strict margins of the preceding construction.
\bibitem{sourceqdn}
E.~Sussman and collaborators,
\href{https://github.com/ethanwsussman/Yang--Mills-LLM-experiment}
{Yang--Mills-LLM-experiment}, deposited raw accounts at commit
\texttt{e392ef45013a581664bab76b90867170b30fe10e}.
\end{thebibliography}
\end{document}
